El Sistema de Control de Comparendos Policiales (SCMC) del municipio de Palermo, un proyecto académico desarrollado en C\# .NET 8 y Angular 16+, se encuentra en una fase avanzada de definición de su lógica de negocio, habiendo estabilizado sus módulos de gestión de infracciones y acuerdos de pago. Sin embargo, la documentación de requisitos del sistema aún contempla un proceso de recaudo manual y presencial, el cual representa un cuello de botella operativo. Este artículo propone y justifica la integración de la pasarela de pago Wompi como la evolución técnica necesaria. Basado en el análisis de casos de estudio análogos (Articulo 10), se establece que el principal riesgo operacional en sistemas de recaudo es la aplicación manual diferida de pagos (24-48 horas). Proponemos una arquitectura de Integración Servidor-a-Servidor (S2S) mediante Webhooks, un modelo que garantiza la conciliación inmediata y la integridad de los datos. Esta implementación se alinea con la estrategia de "time-to-market" (Articulo 2) y es financieramente viable (Articulo 10). El éxito de la adopción ciudadana se supedita a los requisitos de Seguridad Percibida, exigiendo un frontend (Angular) con Responsive Web Design (RWD) y la visibilidad del ícono PSE (Articulo 16).